\vspace{-2ex}
\section{Introduction}
\vspace{-1ex}

Blockchain technology has paved the way for decentralized, resilient, and
programmable ledgers on a global scale. One of its most impactful applications
is smart contracts. These smart
contracts allow developers to encode intricate transaction rules that govern
the ledger. This innovation has made both blockchains and smart contracts
essential infrastructure for decentralized financial services,
commonly known as DeFi. As of Sept 25th, 2023, the Total digital asset Value Locked (TVL) in
$2,933$ DeFi protocols has reached an impressive $48.58$
billion~\cite{defillama2}.

However, the landscape is not without its challenges. Security attacks pose a
significant threat to the security of smart contracts. Attackers can exploit
various vulnerabilities by sending malicious transactions, potentially leading
to the theft of millions of dollars. As of Sept 25th, 2023, the financial
losses attributed to security attacks on DeFi protocols exceeded $5.53$ billion
USD~\cite{defillama}.

One key observation is that transactions initiated by attackers often display
abnormal behaviors when compared to standard transactions from regular DeFi
contract users. These malicious transactions may exploit control flows in
corner cases, use abnormally large values to trigger overflows, or manipulate a
large volume of digital assets to distort the market in DeFi contracts. In
fact, industry experts have been actively monitoring abnormal digital asset
movements on-chain to report malicious activities. For example, Forta Network~\cite{Forta}
deploys monitoring bots to detect on-chain security-related events in real-time. 
Driven by this observation, smart contract developers have proposed deploying
runtime checks to detect transactions leading to abnormal behaviors to neutralize
malicious attacks. These checks involve enforcing various runtime invariants,
such as restricting the maximum number of digital asset deposits or withdrawals
in a contract to prevent market manipulation. Another example is to limit the
interaction between contracts to prevent attackers from crafting sophisticated
attack strategies. However, these mechanisms are often manually designed and
tailored for specific DeFi protocols. This raises questions about their
effectiveness across different types of contracts and whether they maintain an
acceptable false positive rate without hindering normal user activities.

\noindent \textbf{Smart Contract Invariant Study:}
This paper presents the first comprehensive, quantitative analysis focused on
the utilization of dynamically inferred invariants to enhance smart contract security. We
examine $23$ invariant templates, which are advocated by leading auditing firms,
academic research, and DeFi protocol developers. Our findings indicate that
dynamically inferred dynamic invariants serve as effective mechanisms for 
thwarting security
breaches. This paper then proposes new strategies to combine multiple invariants effectively.
Our combined invariants can neutralize over
74.1\% of malicious attacks while maintaining a false positive rate of less than
0.28\%.

\noindent \textbf{{\name}:}
To facilitate this study, we have developed {\name}, a scalable and extensible
invariant synthesis framework. {\name} is designed to automatically derive
invariants from transaction traces through the use of trace and dynamic taint analysis. 
{\name} leverages the main feature of public blockchains, 
transparent databases of transactions histories containing well-organized transaction execution data. 
Then, for each invariant template under consideration, {\name} employs a specialized inference
algorithm to dynamically generate the corresponding invariant based on
historical transaction data. 

\noindent \textbf{Experimental Results:}
We evaluate {\name} on a benchmark set of $42$ smart contracts that have
previously fallen victim to security attacks. Our results show that properly
constructed invariants are effective in neutralizing security threats in $39$
out of the $42$ benchmark contracts. 
%Moreover, this efficacy is achieved with an
%average false positive rate of just $2.98$\%.

In the course of our study, we categorized the $23$ invariant templates into
eight distinct groups based on their underlying design principles: access
control, time lock, gas control, re-entrancy, oracle, storage, money flow, 
and data flow. Subsequently, we conducted a series of in-depth analyses to compare
the efficacy of invariants within each group. We also manually scrutinized the
transactions flagged by each invariant template, leading to several key
findings:
\vspace{-1mm}
\begin{itemize}[leftmargin=*]
\item \textbf{Finding 1:} Certain invariant outperform others in
    terms of effectiveness. Within each invariant group, we identified at least one
        invariant template that is quantitatively superior, neutralizing a
        greater number of attacks while generating fewer false positives.
        See Section~\ref{sec:rq1}.

\item \textbf{Finding 2:} Invariants remain effective even when attackers are
    aware of them in the majority of cases. A common concern regarding runtime
        invariants is their potential vulnerability to informed attackers. Our
        study reveals that selected invariants in the access control, time lock,
        gas control,
        money flow, and data flow groups often directly counter critical elements of
        attack strategies, such as flash loans and transaction atomicity. These
        invariants not only neutralize the malicious transactions but also
        render the attack strategies unfeasible or non-profitable in 84.21\% of
        cases. See Section~\ref{sec:rq2}.

\item \textbf{Finding 3:} Normal users can possibly circumvent invariant guards, 
        thereby mitigating the impact on user
        experience. For example, in the case of data flow and money flow
        invariants, a user can divide a large transaction into smaller segments
        to bypass the invariant guard in 80\% false positive instances. 
        See Section~\ref{sec:rq2}.

\item \textbf{Finding 4:} Combined invariants, formed through disjunction or
    conjunction, offer enhanced security coverage with lower false positive rate.
    Different groups of
        invariants address different attack scenarios. A combined invariant
        formed through conjunction can cover more attack vectors, while one
        formed through disjunction may reduce the false positive rate, as
        malicious attacks often exhibit multiple abnormal behaviors. See
        Section~\ref{sec:rq3}.
\end{itemize}
\noindent \textbf{Contributions:} This paper presents the following contributions:
\begin{itemize}[leftmargin=*]
    \item \textbf{Invariant Inference:} This paper conducts an extensive study
        of $23$ invariant templates, categorized into $8$ distinct groups.
        Additionally, we introduce innovative techniques for the effective
        inference of invariants across all studied templates from transaction 
        history. 

    \item \textbf{{\name}:} This paper presents the design and implementation of {\name},
        a specialized tool for smart contract trace analysis that is capable of
        inferring the invariants under study from transaction history.

    \item \textbf{Experimental Results:} This paper presents the first
        systematic and quantitative evaluation of the effectiveness of runtime
        invariants on $42$ victim contracts in $27$ real-world exploits with high
        financial losses.
        
        % Our empirical findings
        % demonstrate that appropriately crafted invariants can effectively guard
        % against $22$ out of $27$ benchmark attacks, while maintaining a
        % low false positive rate of 0.4\%.

    \item \textbf{Invariant Study Findings:} Our research uncovers a series of
        critical insights that will inform the future application and
        development of dynamic invariants.
\end{itemize}

% The remaining of this paper is organized as follows. XXX
